(1) 长度为 \(L\)、抗弯刚度 \(B=EI>0\) 的不可伸长平面细杆，两端无弯矩铰支，受大小为 \(P>0\) 的相向水平端力。用弧长 \(s\) 与切线角 \(\theta(s)\) 描述其第一对称屈曲分支：左端角为 \(\theta_0\in(0,\pi)\)，中点角为零，右端角为 \(-\theta_0\)，端点弯矩为零。忽略剪切变形与自重。由弹性能与端力势建立非线性平衡方程及第一积分。
(2) 令 \(k=\sin(\theta_0/2)\)，以完全椭圆积分 \(K(k)=\int_0^{\pi/2}(1-k^2\sin^2u)^{-1/2}du\) 表示第一屈曲分支的压缩载荷 \(P\)，并求小幅临界载荷。
(3) 求端点沿载荷方向的距离 \(X=\int_0^L\cos\theta\,ds\) 与端缩 \(\Delta=L-X\)，写成 \(K(k)\) 和 \(E(k)=\int_0^{\pi/2}\sqrt{1-k^2\sin^2u}\,du\) 的闭式。
(4) 计算总弯曲能 \(U_B=(B/2)\int_0^L\theta'^2ds\)，检验 \(k\to0\) 的二阶能量与 \(P/P_{\rm cr}\) 的首项，并解释为何线性欧拉方程不足以确定后屈曲端缩。
